% \section{Background}
\section{Background and Motivation}

\subsection{Data Center Growth and Electricity Pricing}
\edit{BY AKSHAY, GRANT TAKE A LOOK}
The recent growth of large scale data center loads has become a focal point in public discourse surrounding electricity pricing and grid upgrades. The determination of retail rates is multi-factorial, influenced by tariffs, fuel prices, regulatory decisions, infrastructure investments, and resource adequacy procurements. Still, the impacts of large load siting decisions in the wholesale market have long reaching effects that may impact consumers. 

The injection of large loads at specific locations on the grid can trigger congestion driven redispatch. By altering power flows, new load can cause transmission constraints to bind, requiring the operator to redispatch generation away from cheaper units to more expensive marginal resources. This can increase system operating costs and raise \emph{locational marginal prices (LMPs)} at nodes affected by the induced congestion. Although LMPs are wholesale settlement prices rather than retail tariffs, they are an indicator of congestion making load more expensive to serve. Other large loads may be exposed directly to these higher nodal energy prices, while consumers can be affected indirectly through rate adjustments that seek to recover higher system costs over time. Large load interconnection is often accompanied by additional infrastructure costs (e.g. transmission expansion) to handle congestion, and these costs may also be socialized among ratepayers. 

Of course, attributing changes in retail rates directly to a specific load interconnection is not possible. However, it is important to understand how siting large data center loads contributes to changing congestion patterns and system-wide pricing signals that may be directionally consistent with increased downstream cost pressure. Addressing this question requires a careful evaluation of the system as a whole, as congestion and redispatch effects can propagate non-locally across the network. 

\subsection{Existing Data Center Planning Practices}
Large load planning on transmission networks typically evaluates system reliability through stressed-case snapshots (e.g., peak and off-peak conditions) to ensure voltage bounds, thermal limits, and stability margins are maintained~\cite{nerc_tpl001_5.1}. This approach treats loads as a capacity requirement and asks whether sufficient headroom exists to serve it reliably, assuming other siting and regulatory features have already been met. Datacenter siting decisions, in turn, prioritize power availability, land cost, fiber connectivity, and natural disaster risk, treating the grid as a constraint to satisfy rather than a system whose behavior changes with placement~\cite{goiri2011intelligentplacement,ALAPERA201869}. While grid-side planners consider some congestion effects, neither process systematically can quickly evaluate how different mechanics of load injection affect congestion patterns, wholesale market pricing, or dispatch costs; these variables depend jointly on where capacity connects, how much is injected, and when demand occurs. Recent NERC guidance acknowledges this gap, noting that the operational behavior of emerging large loads like datacenters can induce new risks that fall outside of the previous ranges; recommending that planners should explore several new strategies and load patterns~\cite{nerc2026}. It has been well documented that AI workloads prevalent in these new large loads exhibit fast, large variability across training and inference phases~\cite{chen2025electricitydemandgridimpacts, patel2024llms, li2024unseenaidisruptionspower, dynamollm}. This variance implies that the same provisioned MW can produce materially different grid congestion and pricing effects depending on workload mix. Our framework addresses these limitations by coupling temporal load profiles with planning and dispatch levers to allow for an evaluation of placement strategies under these different operational conditions.

\subsection{Mitigation Levers}

There are a plethora of conventional planning interventions to mitigate the impacts of connecting large data center loads—building or upgrading transmission infrastructure, expanding generation resources, or curtailing loads. These levers are invaluable tools in managing the health of the grid, but can also be slow and capital intensive. Additionally, many of these levers are positioned as \emph{system} upgrades in \emph{response} to an interconnection request, rather than as a design variable on the load side. 

We posit that the \emph{distribution} of large data center loads offers a load-side lever that can be used in conjunction with more traditional planning levers. By choosing where to concentrate loads on the grid, siting decisions act as a way to reduce externalities such as peak-price events without requiring immediate system-wide upgrades or reducing the magnitude of such upgrades. Distribution can reduce the burden of congestion costs and redispatch impacts, which may translate into lower wholesale exposure for the data center operator and lower long-term cost pressure for the system as a whole. 

This additional planning tool comes with important caveats. Distributed placement can increase a data center operator's operational complexity, with additional care required to handle latency constraints, supply-chain logistics, and land acquisition. Load curtailment and flexible operation are also possible options to mitigate the effects of data center loads, but deploying these options reliably at a large-scale requires additional study which is out of the scope of this work. 

We emphasize that distribution as a load-side lever is meant to be used alongside other planning levers. Distributed interconnection paired with transmission expansion or storage deployment may reduce the magnitude of upgrades required to service the same aggregate demand. Considering load placement as complementary to other grid invests can ``shift the knee" in price response curves, thereby expanding the range of demand levels that can be served without severe price spiking. 

% \subsection{Capacity Expansion Modeling}
% Capacity expansion problems, and more specifically bilevel capacity expansion problems, have a rich literature, particularly from the perspective of planning grid infrastructure investments. A well-studied class of problems closely related to our setting is \textit{profit maximization}, where merchant investors plan new generation investments with the goal of maximizing market profits while accounting for the need to upgrade transmission infrastructure when necessary. In these models, investors seek to favorably influence market outcomes while minimizing expenditures on costly grid reinforcements \cite{baringo2011wind, wogrin2011generation}.

% These problems are naturally thought of as bilevel optimization problems and can be viewed as a type of Stackelberg game, where a leader (the investor) anticipates the market clearing outcome (the follower). In general, bilevel expansion problems are highly nonconvex and computationally challenging. Classical approaches aim to find globally optimal solutions by first reformulating the problem as a Mathematical Program with Equilibrium Constraints (MPEC). These MPECs can then be further reformulated as mixed-integer programs (MIPs) that are solvable with commercial solvers such as Gurobi \cite{gurobi2025}. While this approach theoretically enables exact solutions, it comes at a significant computational cost. Bilevel expansion problems are known to be NP-hard, and solving large MIP formulations to global optimality quickly becomes infeasible, particularly when the planner wishes to evaluate a large number of scenarios with varying system conditions and uncertainty realizations.

% In this work, we use the gradient-based methods for bilevel expansion planning introduced in \cite{degleris2024gradient}, which aim to find high-quality, locally optimal solutions much more efficiently than traditional MIP-based approaches. Gradient methods avoid the combinatorial explosion associated with MPEC reformulation and instead directly leverage the structure of the bilevel problem to enable scalable optimization. Given the complex, coupled constraints associated with data center buildout—including grid infrastructure limitations, quality-of-service (QoS) requirements, and capital expenditure (CapEx) considerations—it is critical that the planner can quickly iterate over many candidate scenarios to identify acceptable and actionable planning solutions.


% \subsection{Data Center Power and Siting Constraints}
% Prior work has extensively studied power modeling and planning for data centers, though typically without considering the broader grid context that our approach addresses. Power distribution hierarchies in data centers involve multiple distribution paths for high availability and fault tolerance~\cite{hsu2018smoothoperator,zhang2021flex, wu2016dynamo}. Investment decisions in placement of resources are therefore made at a rack level due to the capital expenditures required for the surrounding infrastructure. Beyond the actual hardware, workload characterization like in Coach~\cite{reidys2025coach}, DynamoLLM~\cite{dynamollm}, and POLCA~\cite{patel2024llms}, demonstrated that general data center and AI inference workloads exhibit temporal patterns with complementary usage across different applications, making them suitable for strategic placement and oversubscription. Although these works provide valuable information on internal data center power management and workload characterization, they do not address the cross-layer optimization between data center placement and grid capacity expansion that our approach aims to achieve.

% \subsection{Cost Concerns for Large Datacenter Loads}
% A recent discussion in news and government has concerned the harmful
% effects of data centers on local communities, namely the rising cost
% in electric utility bills. While there are many confounding variables,
% adding large data center loads to areas does have the potential to
% drive up costs for residential and other commercial consumers, as
% often grid upgrades are socialized 
